\section{What this software is}

This repository is the digital twin of the pulse-echo azimuthal
measurement: a rotational single-transducer ultrasound method that
recovers the crystal orientation fabric (COF) of a polycrystalline ice
disk from the azimuthal variation of diametral sound velocity. The
physical instrument is documented in its own repository and manual;
this codebase answers the questions an instrument cannot: what should
the measurement see for a known fabric, which parts of the received
signal carry fabric information, how do acquisition errors corrupt it,
and how well can the fabric be recovered.

The pipeline has four stages, mirrored by the GUI tabs:

\begin{enumerate}
  \item \textbf{Specimen.} Build a synthetic disk: Voronoi grain
    structure with controlled size distribution, and a Watson
    orientation distribution with controlled strength and mean
    direction (Chapter~\ref{ch:specimen}).
  \item \textbf{Forward model.} Propagate pulses through the disk,
    either with the fast ray model (interactive) or the label-indexed
    anisotropic FDTD solver (full-wave, GPU), which is itself
    validated against an independent finite-element solver
    (Chapter~\ref{ch:forward}).
  \item \textbf{Acquisition and errors.} Simulate the rotational
    acquisition of the physical machine, with its error model
    (Chapter~\ref{ch:acquisition}).
  \item \textbf{Analysis and inversion.} Recover the fabric from the
    simulated traces and quantify what is and is not recoverable
    (Chapter~\ref{ch:analysis}).
\end{enumerate}

The repository doubles as the evidence base for the journal
manuscript on the method. The working convention throughout: every
number quoted in the paper traces to a module here, and the expected
output is recorded in that module's docstring. Negative results are
kept deliberately, because ruling a signal path out is as much a part
of the argument as ruling one in.
